% chapters/03/section-05.tex
\section{Program Pygame Pertama}
\label{sec:program-pertama}

Bagian ini menuntun Anda menulis program Pygame yang benar-benar berjalan.
Sasaran programnya sederhana: membuka jendela, menampilkan informasi versi
perkakas yang dipakai, dan menunggu sampai jendela ditutup.

\subsection{Kerangka Program dalam Lima Langkah}

Setiap program Pygame, sederhana maupun rumit, mengikuti kerangka yang sama.
Kelima langkah berikut akan Anda temui kembali, dengan tambahan, pada seluruh
bab berikutnya.

\begin{langkah}
  \item \textbf{Memulai Pygame} dengan \code{pygame.init()}, sehingga seluruh
        bagian pustaka yang diperlukan siap dipakai.
  \item \textbf{Membuat permukaan tampilan} dengan
        \code{pygame.display.set_mode()}, yang membuka jendela dan
        mengembalikan permukaan tempat Anda menggambar.
  \item \textbf{Menggambar isi jendela}: mengisi latar dengan satu warna, lalu
        menempelkan teks atau bentuk di atasnya.
  \item \textbf{Menayangkan hasil gambar} dengan \code{pygame.display.flip()},
        karena gambar yang telah disusun masih berada di memori dan belum
        terlihat pemain sebelum perintah itu dijalankan.
  \item \textbf{Menunggu sampai program diakhiri, lalu membersihkan} dengan
        \code{pygame.quit()}.
\end{langkah}

\subsection{Kode Lengkap}

\kodepython{python/ch03_jendela.py}{Program Pygame pertama: jendela dan informasi versi}{lst:ch03-jendela}

\subsection{Membedah Kode}

\subsubsection{Membaca pilihan dari baris perintah}

Fungsi \code{parse_arguments} memakai modul \code{argparse} bawaan Python untuk
membaca pilihan yang Anda berikan ketika menjalankan program. Karena itu,
ukuran jendela dapat diubah tanpa menyunting kode. Cara ini dipakai pada
hampir seluruh program di buku ini, sebab mengubah nilai melalui baris
perintah jauh lebih praktis daripada mengubah kode dan menjalankannya
berulang kali.

Pilihan \code{--auto-close} disediakan untuk pengujian otomatis. Ketika
pilihan itu dipakai, program mengirim sendiri peristiwa penutupan jendela
sehingga program berakhir tanpa perlu disentuh. Pilihan semacam ini berguna
ketika program harus diuji pada komputer tanpa layar.

\subsubsection{Memulai Pygame dan membuka jendela}

Pemanggilan \code{pygame.init()} menyalakan seluruh bagian Pygame: tampilan,
huruf, bunyi, dan penanganan peristiwa. Sesudah itu,
\code{pygame.display.set_mode((640, 480))} membuka jendela berukuran
640 kali 480 piksel. Nilai yang dikembalikannya disimpan pada variabel
\code{screen}, yaitu \textbf{permukaan} (\istilah{surface}) tempat seluruh
penggambaran dilakukan.

Istilah permukaan perlu Anda pahami sejak awal. Pada Pygame, apa pun yang
digambar---jendela, gambar yang dimuat dari berkas, atau teks---adalah
permukaan. Menggambar berarti menempelkan satu permukaan ke permukaan lain
dengan perintah \code{blit}. Cara berpikir itu akan terasa wajar setelah Anda
memakainya beberapa kali.

\subsubsection{Warna, huruf, dan teks}

Warna dinyatakan sebagai rangkap tiga bilangan antara 0 dan 255 yang
menyatakan banyaknya merah, hijau, dan biru. Nilai \code{(31, 56, 100)}
menghasilkan biru tua, sedangkan \code{(255, 255, 255)} menghasilkan putih.

Huruf dibuat dengan \code{pygame.font.Font(None, 26)}. Argumen pertama bernilai
\code{None}, yang berarti Pygame memakai huruf bawaan. Pilihan itu disengaja
agar program tidak bergantung pada huruf yang terpasang pada sistem operasi
tertentu. Argumen kedua adalah ukuran huruf dalam piksel.

Teks belum menjadi gambar sampai dijalankan melalui \code{font.render()}.
Fungsi itu mengubah rangkaian huruf menjadi permukaan berisi gambar teks,
beserta keterangan bahwa latar belakangnya dibuat transparan. Permukaan hasil
\code{render} itulah yang kemudian ditempelkan ke jendela dengan \code{blit}.

\subsubsection{Menunggu peristiwa dan mengakhiri program}

Perhatikan bagian berikut dengan saksama, karena di sinilah letak perbedaan
antara program ini dan siklus permainan yang sebenarnya.

\begin{kodekeluaran}
running = True
while running:
    event = pygame.event.wait()
    if event.type == pygame.QUIT:
        running = False
\end{kodekeluaran}

Perintah \code{pygame.event.wait()} menghentikan program sampai satu peristiwa
(\istilah{event}) terjadi, misalnya tombol ditekan atau jendela ditutup.
Apabila peristiwa itu adalah \code{QUIT}, yaitu permintaan menutup jendela,
variabel \code{running} diubah menjadi salah dan perulangan berakhir.

Perulangan tersebut adalah \textbf{perulangan peristiwa}, bukan siklus
permainan. Ia hanya menunggu dan bereaksi; tidak ada keadaan yang diperbarui
dan tidak ada gambar yang disusun ulang. Akibatnya, apa pun yang terjadi pada
jendela---misalnya ditutup sebagian oleh jendela lain---isi jendela tidak
digambar kembali. Untuk permainan yang sesungguhnya, ketiga pekerjaan itu,
yaitu membaca masukan, memperbarui keadaan, dan menggambar ulang, harus
dilakukan berulang kali tanpa berhenti. Itulah yang dibahas pada Bab~IV.

Pemanggilan \code{pygame.quit()} pada akhirnya mematikan kembali seluruh
bagian Pygame. Kebiasaan memanggilnya sebelum program berakhir mencegah
berbagai masalah, terutama ketika program dijalankan berulang kali dari
editor yang sama.

\subsection{Percobaan yang Disarankan}

Cara tercepat memahami sebuah program bukan dengan membacanya berulang kali,
melainkan dengan mengubahnya. Ubahlah satu hal pada satu waktu, jalankan
kembali, dan amati bedanya. Tabel~\ref{tab:percobaan-ch03} memuat percobaan
yang dapat Anda mulai.

\begin{table}[htbp]
\centering
\caption{Percobaan sederhana pada program pertama}
\label{tab:percobaan-ch03}
{\setlength{\tabcolsep}{4pt}\footnotesize%
\begin{tabularx}{\textwidth}{@{}L{6.0cm} Y@{}}
\toprule
\headrow \textbf{Yang diubah dan caranya} & \textbf{Yang sebaiknya Anda amati} \\
\midrule
Ukuran jendela, melalui pilihan \code{--width} dan \code{--height} pada baris
perintah. &
Letak teks tetap pada sudut kiri-atas, sedangkan ruang kosong di sebelah kanan
dan bawah bertambah. \\
\addlinespace
Warna latar, melalui nilai \code{WARNA_LATAR}. &
Perubahan pada ketiga angka mengubah warna secara berbeda. Naikkan angka
pertama untuk memperkuat merah. \\
\addlinespace
Ukuran huruf, melalui angka 26 pada pembuatan huruf. &
Angka yang lebih besar menghasilkan teks yang lebih tinggi. Tambahkan baris
teks untuk melihat jarak antar-baris. \\
\addlinespace
Jarak antar-baris, melalui nilai 32 pada penambahan posisi tegak. &
Nilai yang lebih kecil membuat baris teks saling menempel atau saling
menimpa. \\
\addlinespace
Judul jendela, melalui teks pada \code{pygame.display.set_caption}. &
Judul pada bilah jendela berubah. Judul ini berguna untuk menandai program
mana yang sedang berjalan. \\
\bottomrule
\end{tabularx}}
\end{table}
